\documentclass[10pt,a4paper]{amsart}
\usepackage[margin=19mm]{geometry}
\usepackage{amsmath,amssymb,mathtools}
\usepackage[scr=rsfs,cal=euler]{mathalfa}
\usepackage{xfrac}
\usepackage[unicode,hidelinks,bookmarks=false]{hyperref}
\newcommand{\ie}{i.e.\ }
\newcommand{\cf}{cf.\ }
\newcommand{\resp}{respectively\ }
\newcommand{\Subsection}{\subsection}
\providecommand{\nobreakdash}{\nobreak\hbox{-}}
\setlength{\emergencystretch}{2em}
\multlinegap0pt
\usepackage{bm}
\usepackage{bbm}
\usepackage{dsfont}
\usepackage{xspace}
\usepackage{cancel}
\usepackage{mathtools}
\usepackage{amsthm}
\makeatletter
\@ifundefined{theorem}{\newtheorem{theorem}{Theorem}}{}
\@ifundefined{lemma}{\newtheorem{lemma}[theorem]{Lemma}}{}
\makeatother
\title[Exact Recovery in the Data Block Model]{Exact Recovery in the Data Block Model}
\author{Amir R. Asadi, Akbar Davoodi, Ramin Javadi,, Farzad Parvaresh}
\date{Source: arXiv:2602.05852v1 (CC BY 4.0)}
\begin{document}
\maketitle

\noindent\emph{Source and licence.} Exact Recovery in the Data Block Model. Version arXiv:2602.05852v1 (CC BY 4.0), distributed under the Creative Commons Attribution 4.0 International licence, as recorded in the arXiv licence metadata for this exact version and shown on the abstract page https://arxiv.org/abs/2602.05852v1. Full rights notice: licenses/arxiv-2602.05852-CC-BY-4.0.txt.\par\smallskip

\noindent\textit{Editorial scope.} Counted unit: the Lemma whose source label is \texttt{thm:scheffe} in the article (source0.tex lines 799--837, source label \texttt{thm:scheffe}), delivered together with the article's own complete original proof of it (source0.tex lines 838--945, one \texttt{proof} environment of 108 lines). No mathematics is written, repaired or re-derived: statement and proof are the article's own text line for line. Required context retained uncounted: none. Every definition, display and label of those lines is retained unchanged. Omitted from this package and not redistributed: 1--798, 946--1645. Nothing inside a retained excerpt is omitted or altered: every retained body is delivered exactly as the article's lines are written, so no omission receipt of any kind accompanies this package. Citations: the retained excerpts carry no citation command at all, so no bibliography entry of the article is transcribed and no reference is redistributed. Reference adaptation: none was needed and none was made; every reference inside the delivered unit resolves inside it, so no unresolved reference or citation is produced. Printed slips kept literal: no printed word, symbol, hypothesis or number of the counted statement or of its proof was altered. Layout and class: the article's own class is \texttt{article} with packages including geometry, amsmath, amssymb, amsfonts, amsthm, bm, bbm, mathrsfs, graphicx, grffile, booktabs, subcaption, authblk, algorithm; the wrapper is the lane's pinned \texttt{amsart} editorial wrapper at 10pt on A4, which restates the theorem-like environments the retained text needs and adds the article's own macro definitions that the retained text needs (none), with extra wrapper packages (bm, bbm, dsfont, xspace, cancel, mathtools). The delivered numbering is the unit's own. Assets: no retained span contains a figure environment, a graphics-inclusion command, a PDF-page inclusion, a table or a TikZ picture; no image file is copied into this package, the package declares no asset dependency and no image file is redistributed with it. Licence: version arXiv:2602.05852v1 of this article is distributed under the Creative Commons Attribution 4.0 International licence, as recorded in the arXiv licence metadata for this exact version and shown on the abstract page https://arxiv.org/abs/2602.05852v1. A full rights notice is delivered at licenses/arxiv-2602.05852-CC-BY-4.0.txt. Characters: every retained excerpt is pure ASCII, so the delivered engine, XeLaTeX with its default Latin Modern text font, reports no missing character.
\begin{lemma}[Scheff\'e test under $m$ corruptions]\label{thm:scheffe}
Let $P_X^{(1)}$ and $P_X^{(2)}$ be probability mass functions on a discrete alphabet $\mathcal{X}$.
Let $X_1,\dots,X_n$ be observations such that there exists an (unknown) index set
$\mathcal{I}\subseteq[n]$ with $|\mathcal{I}|\ge n-m$ for which the random variables
$\{X_i:i\in\mathcal{I}\}$ are i.i.d.\ with common law $P_X$, where
\[
P_X\in\{P_X^{(1)},P_X^{(2)}\},
\]
while the remaining $m$ observations $\{X_i:i\notin\mathcal{I}\}$ are arbitrary.

Consider testing
\[
H_1:\ P_X=P_X^{(1)}
\qquad\text{vs.}\qquad
H_2:\ P_X=P_X^{(2)}.
\]
Let the empirical measure be
\[
P_n(A):=\frac{1}{n}\sum_{i=1}^n \mathbbm{1}\{X_i\in A\},\qquad A\subseteq\mathcal{X},
\]
and define
\[
A^*=\arg\max_{A\subseteq\mathcal{X}} \bigl|P_X^{(1)}(A)-P_X^{(2)}(A)\bigr|.
\]
Consider the Scheff\'e decision rule
\[
\bigl|P_n(A^*)-P_X^{(1)}(A^*)\bigr|
\underset{H_1}{\overset{H_2}{\gtrless}}
\bigl|P_n(A^*)-P_X^{(2)}(A^*)\bigr|.
\]
If
\[
\mathrm{TV}\!\left(P_X^{(1)},P_X^{(2)}\right)>\frac{2m}{n},
\]
then the (prior-averaged) Bayes error probability $P_e$ of this test satisfies
\[
P_e \le 4\exp\!\left(-\frac{n}{2}\Bigl(\mathrm{TV}\!\left(P_X^{(1)},P_X^{(2)}\right)-\frac{2m}{n}\Bigr)^2\right).
\]
\end{lemma}

\begin{proof}
Let $A^*\in\arg\max_{A\subseteq\mathcal X}\bigl|P_X^{(1)}(A)-P_X^{(2)}(A)\bigr|$, so that
\[
\bigl|P_X^{(1)}(A^*)-P_X^{(2)}(A^*)\bigr|=\mathrm{TV}\!\left(P_X^{(1)},P_X^{(2)}\right).
\]
Define $S_n:=\sum_{i=1}^n \mathbbm{1}\{X_i\in A^*\}$, so that $P_n(A^*)=S_n/n$.

We first bound the type-I error under $H_1$ (i.e., $P_X=P_X^{(1)}$ on the uncorrupted samples).
An error occurs only if the empirical mass $P_n(A^*)$ is closer to $P_X^{(2)}(A^*)$ than to $P_X^{(1)}(A^*)$, i.e.
\[
\bigl|P_n(A^*)-P_X^{(1)}(A^*)\bigr|>\bigl|P_n(A^*)-P_X^{(2)}(A^*)\bigr|.
\]
By the triangle inequality,
\[
\bigl|P_X^{(2)}(A^*)-P_X^{(1)}(A^*)\bigr|
\le
\bigl|P_n(A^*)-P_X^{(1)}(A^*)\bigr|
+
\bigl|P_n(A^*)-P_X^{(2)}(A^*)\bigr|,
\]
thus on the error event we must have
\[
2\bigl|P_n(A^*)-P_X^{(1)}(A^*)\bigr|
>
\bigl|P_X^{(2)}(A^*)-P_X^{(1)}(A^*)\bigr|
=
\mathrm{TV}\!\left(P_X^{(1)},P_X^{(2)}\right).
\]
Hence
\begin{equation}\label{eq:typeI_reduce}
P_e(1)
\le
\Pr_{H_1}\!\left\{\,2\bigl|P_n(A^*)-P_X^{(1)}(A^*)\bigr|
>
\mathrm{TV}\!\left(P_X^{(1)},P_X^{(2)}\right)\right\}.
\end{equation}

Let $\mathcal I\subseteq[n]$ be the (unknown) index set of uncorrupted samples, with $|\mathcal I|\ge n-m$.
Write
\[
S_n = \sum_{i\in\mathcal I}\mathbbm{1}\{X_i\in A^*\} \;+\; \sum_{i\notin\mathcal I}\mathbbm{1}\{X_i\in A^*\}
=: S_{\mathcal I}+S_{\mathcal I^c}.
\]
Under $H_1$, the variables $\{\mathbbm{1}\{X_i\in A^*\}: i\in\mathcal I\}$ are i.i.d.\ Bernoulli with mean
$P_X^{(1)}(A^*)$, while $S_{\mathcal I^c}\in[0,m]$ is arbitrary.
Therefore,
\[
\bigl|S_n-(n-m)P_X^{(1)}(A^*)\bigr|
\le
\bigl|S_{\mathcal I}-(n-m)P_X^{(1)}(A^*)\bigr| + m,
\]
and dividing by $n$ gives
\[
\bigl|P_n(A^*)-P_X^{(1)}(A^*)\bigr|
=
\frac{1}{n}\bigl|S_n-nP_X^{(1)}(A^*)\bigr|
\le
\frac{1}{n}\bigl|S_{\mathcal I}-(n-m)P_X^{(1)}(A^*)\bigr|+\frac{m}{n}.
\]
Plugging this into \eqref{eq:typeI_reduce} yields
\[
P_e(1)
\le
\Pr_{H_1}\!\left\{
\frac{2}{n}\bigl|S_{\mathcal I}-(n-m)P_X^{(1)}(A^*)\bigr|
>
\mathrm{TV}\!\left(P_X^{(1)},P_X^{(2)}\right)-\frac{2m}{n}
\right\}.
\]
By Hoeffding's inequality for sums of $(n-m)$ i.i.d.\ Bernoulli variables,
\begin{align*}
	&\Pr_{H_1}\!\left\{
\bigl|S_{\mathcal I}-(n-m)P_X^{(1)}(A^*)\bigr|
>
\frac{n}{2}\Bigl(\mathrm{TV}(P_X^{(1)},P_X^{(2)})-\frac{2m}{n}\Bigr)
\right\}\\
&\quad\quad\quad\le
2\exp\!\left(
-\frac{1}{2}n\Bigl(\mathrm{TV}(P_X^{(1)},P_X^{(2)})-\frac{2m}{n}\Bigr)^2
\right),
\end{align*}
where we used $n-m\le n$ to simplify the exponent.
Hence,
\[
P_e(1)\le
2\exp\!\left(
-\frac{1}{2}n\Bigl(\mathrm{TV}(P_X^{(1)},P_X^{(2)})-\frac{2m}{n}\Bigr)^2
\right).
\]
The same argument under $H_2$ gives
\[
P_e(2)\le
2\exp\!\left(
-\frac{1}{2}n\Bigl(\mathrm{TV}(P_X^{(1)},P_X^{(2)})-\frac{2m}{n}\Bigr)^2
\right).
\]
Finally, for any prior on $\{H_1,H_2\}$, the Bayes error satisfies
\(
P_e \le P_e(1)+P_e(2),
\)
and therefore
\[
P_e \le
4\exp\!\left(
-\frac{1}{2}n\Bigl(\mathrm{TV}(P_X^{(1)},P_X^{(2)})-\frac{2m}{n}\Bigr)^2
\right).
\]
\end{proof}

\end{document}
